\section{奖励从哪里来}

在之前的讲次中，我们假设奖励函数 $r(s, a)$ 是给定的。但在现实中，\textbf{定义奖励函数本身就是一个难题}。

\begin{knowledgebox}{不同领域的奖励情况}
\begin{itemize}
    \item \textbf{电子游戏}：分数天然就是奖励，直接可用。
    \item \textbf{机器人}：什么是"好的抓取"？什么是"整洁的折叠"？需要人工定义，往往使用代理指标。
    \item \textbf{对话系统}：什么是"好的回答"？高度主观，难以用简单函数衡量。
    \item \textbf{自动驾驶}：涉及安全、舒适、效率等多个维度的权衡。
\end{itemize}
\end{knowledgebox}

除了直接定义奖励外，我们已经见过一种替代方案：\textbf{模仿学习}——直接模仿专家动作，无需奖励。但模仿学习不推理结果和动力学，且专家可能有不同的 embodiment。

\subsection{本章小结}

奖励设计是 RL 应用的关键瓶颈。本讲介绍三种从人类监督中\textbf{学习}奖励的方法。

